\scanpage{118}
1) Les cycles num. équivalents à 0 sur $X$, mod équiv. rationnelle, sont décrits par un nb fini de ``familles algébriques'' de cycles.
\[
\Longrightarrow\left\{
\begin{array}{l}
(\text{équiv. num.}=\tau\text{-équivalence})\quad\Downarrow\\
\text{les }A_a^i(X)\text{ sont de type fini}\ldots
\end{array}\right.
\]
2) $(x,y)-(x,b)-(a,y)+(a,b)\sim_a0$

$1),\text{ et }2)\Longrightarrow$ Les $A_r^i(X)^a$ sont des V.A.
\[
\left\{
\begin{array}{l}
A_r^i(X)^a\times A_r^{\unsure{j|d}}(\unsure{Y})^a\ \unsure{\text{---}|\longrightarrow}\ A_r^{i+\unsure{j|d}}(X\times Y)^a\quad\text{est nul}\\
A_r^i(X)^a\times A_r^{\unsure{j|d}}(X)^a\ \unsure{\text{---}|\longrightarrow}\ A_r^{i+\unsure{j|d}}(X)^a\quad\text{est nul}
\end{array}\right.
\]
\deleted{\unsure{1)}}
\[
A_0^{\unsure{r|*}}(X)^a\xrightarrow{\sim}\operatorname{Alb}(X)\qquad\text{bijectif}
\]
\[
2)\Longleftrightarrow\qquad A_0(X\times Y)^a\ \overset{\sim}{\unsure{\text{---}|\longrightarrow}}\ A_0(X)^a\times A_0(Y)^a\qquad\text{bijectif}
\]
\margin{117}
